
% PUT this before day 14 (i.e., do feature selection before assoc analysis)

%Reducing the feature space
%
%The curse of dimensionality, among other things, makes us want to use fewer
%features than we have.
%
%Question: which ones should we use?
%
%Two basic approaches: feature selection, and dimensionality reduction.
%
%Feature selection:
%
%    Subset selection -- find and use a subset of the features "as is".
%
%        - "best subset": try all possible subsets of features. Computationally
%        intractable.
%
%        - "forward/backward selection": greedy algorithm starting with none,
%        or all, and adding/removing best-looking one at each phase.
%
%    (Mostly for linear regression): Shrinkage/regularization -- use all the
%    predictors, but shrink some towards zero (possibly all the way). Ridge and
%    lasso.
%
%
%Evaluating goodness-of-fit: 
%
%    1) "Mallows' Cp", AIC, BIC, adjusted R-squared (remember: usual R^2 is
%    1-RSS/TSS.) Various ways of penalizing a more complicated model.
%
%    2) Use a training/test split.
%
%
%
%Dimensionality reduction:
%
%    Need linear algebra. "Principal Components" is a common technique, which
%    transforms original features into a new set of orthogonal features, and
%    orders them in decreasing order of how much variance is explained.
%    ("projection")
%
%
%Warning: "the percentage of variance explained"s do not add up to 100\%.
%
%

\documentclass[12pt]{article}

\addtolength{\textwidth}{.8in}
\addtolength{\hoffset}{-.4in}
\addtolength{\textheight}{1in}
\addtolength{\voffset}{-1in}

\usepackage{amsmath}
\usepackage{graphicx}
\usepackage{upquote} % to convert funny quotes to straight quotes
\usepackage{amsfonts}
\usepackage{multicol}
\usepackage{paralist}
\usepackage{icomma}  % no spaces after commas
\setlength{\parindent}{15pt}

\renewcommand*{\thefootnote}{\fnsymbol{footnote}}


\begin{document}

\begin{center}
\begin{Large}
DATA 219 Notes -- Apr. 17, 2018
\end{Large}
\end{center}

\section{Feature selection}

The ``curse of dimensionality" makes us inclined to actually use fewer
features in our classifier than we have at our disposal. This is for two
reasons:

\begin{itemize}
\itemsep.1em
\item Many of the available features are likely to be ``noisy," meaning they
are only loosely correlated with the target feature, if at all. Including
noisy features will make it more difficult for any classifier to separate the
wheat from the chaff: the features which truly are predictive will get drowned
out by those that merely clutter the picture.

\item Even for those features which do contain substantive, reliable
information about the target, it's not likely that they're all
\textit{independent} of each other. If (say) five of our twenty features each
furnished some unique and meaningful information about what the target label
is likely to be, that sounds great...but it's more likely that some (or even
all) of these five are correlated with each other, and hence contain
overlapping information. We'll get a better-performing classifier if we
avoid giving it redundant features that contain pesky noise to boot.

\end{itemize}

So how do we reduce the feature space? There are essentially three approaches:

\begin{enumerate}
\itemsep.1em
\item \textbf{Subset selection}. We'll simply choose a subset of the original
features and use those ``as is" (after normalizing them or one-hot-encoding
them, of course).
\item \textbf{Shrinkage} (a.k.a. ``\textbf{regularization}"). Especially
popular with linear models, we can essentially ``tone down" some of the less
important features by shrinking their coefficients towards zero.
\item \textbf{Transforming the feature space}. In this case, we actually don't
use the \textit{original} features at all; instead, we form a set of
new features out of parts of the original features. This may sound weird --
especially since our new features don't have any obvious ``meaning"; they're
just combinations of some of the original ones -- but it can be very powerful
and informationally efficient.
\end{enumerate}

I'll say much more about the first of these, since the latter two require some
more advanced math.

\subsection{Subset selection}

Let's say we have four features available in our original data set (call them
$F_0,F_1,F_2, \textrm{and} F_3$, to stick with Python's numbering
convention.) As stated above, we may not want to actually use all four. So
which ones could we use?

It turns out we have $2^4$ (which is 16) choices of different feature sets to
use. They are:

\begin{compactitem}
\item all four
\item just $F_0,F_1$ and $F_2$
\item just $F_0,F_1$ and $F_3$
\item just $F_0$ and $F_1$
\item just $F_0,F_2$ and $F_3$
\item just $F_0$ and $F_2$
\item just $F_0$ and $F_3$
\item just $F_0$
\item just $F_1,F_2$ and $F_3$
\item just $F_1$ and $F_3$
\item just $F_1$ and $F_2$
\item just $F_1$
\item just $F_2$ and $F_3$
\item just $F_2$
\item just $F_3$
\item none of them (just use the prior)
\end{compactitem}

The way we get the number ``$2^4$" is to reason as follows: ``each feature can
either be included, or excluded, in any feature set I choose. So for each
of the four features, there are two independent choices (either take it or
leave it) which leaves me with $2 \cdot 2 \cdot 2 \cdot 2$ different possible
sets of features."

The one using all of the features is most likely to overfit, and the one using
none of them is most likely to underfit. Often the best-fitting classifier (on
new data! not the original training data) will be one of the intermediate
ones.

So...how do we find that best-fitting classifier then?

\subsubsection{Consider all subsets}

The obvious solution is to say ``we've identified all possible subsets of
features, let's try 'em!" And indeed we can do this if the number of features
is modest. Train a classifier (whether Naive Bayes, kNN, or anything else)
using only the features on a given set, compute the error rate, and then
choose the set of features with the lowest error rate.

Here's some code that does just that (for Naive Bayes)\footnote{Notes: the
\texttt{is\_numeric\_dtype()} function distinguishes between numeric and
categorical variables, so we can normalize one and one-hot the other. The
\texttt{LabelBinarizer}, as we've seen is one way to one-hot variables;
another approach is the \texttt{.get\_dummies()} method on Pandas DataFrame.
Finally, the ``\texttt{operator}" thing is just a short way of telling Python
``please sort this dictionary, but by value, not by key," so we can get the
results back in percent-accuracy-order rather than
alphabetical-order-of-feature-set.}:

\pagebreak
\begin{center}
\line(1,0){400}
\end{center}
\scriptsize
\begin{Verbatim}[fontsize=\small,samepage=true,frame=none]
import numpy as np
import pandas as pd
from sklearn.preprocessing import LabelBinarizer
from sklearn.naive_bayes import MultinomialNB
from pandas.api.types import is_numeric_dtype

w = pd.read_csv("warriors_extra.csv")

train = w.sample(frac=.7)
test = w.drop(train.index)

# Given a DataFrame and a set of column headings to include, return a NumPy
# matrix with all these features prepared and concatenated. "Prepared" means:
# (1) for categorical variables, one-hot them, and (2) for numeric variables,
# compute their Z-scores.
def prepare_this_puppy(df, include):
    lb = LabelBinarizer()
    prepared_features = []
    for feature in include:
        if is_numeric_dtype(df[feature].dtype):
            thing = df[feature]-df[feature].mean()/df[feature].std()
            prepared_features.append(np.array([thing.as_matrix()]).transpose())
        else:
            prepared_features.append(lb.fit_transform(df[feature]))
    return np.concatenate(prepared_features, axis=1)

feature_sets = [
    {'lightsaber','beard','height','demeanor'},
    {'lightsaber','beard','height'},
    {'lightsaber','beard','demeanor'},
    {'lightsaber','beard'},
    {'lightsaber','height','demeanor'},
    {'lightsaber','height'},
    {'lightsaber','demeanor'},
    {'lightsaber'},
    {'beard','height','demeanor'},
    {'beard','height'},
    {'beard','demeanor'},
    {'beard'},
    {'height','demeanor'},
    {'height'},
    {'demeanor'}]
scores = {}

for feature_set in feature_sets:
    train_features = prepare_this_puppy(train,feature_set)
    test_features = prepare_this_puppy(test,feature_set)

    naive_bayes = MultinomialNB()
    naive_bayes.fit(train_features,train.type.as_matrix())
    nb_predictions = naive_bayes.predict(test_features)
    scores[tuple(feature_set)] = round(
        sum(nb_predictions == test.type)/len(test)*100,2)

import operator
sorted_scores = sorted(scores.items(), key=operator.itemgetter(1), reverse=True)
for feature_set, score in sorted_scores[:10]:
    print("{}: {:.2f}".format(','.join(sorted(feature_set)),score))
\end{Verbatim}
\normalsize
\begin{center}
\line(1,0){400}
\end{center}

\pagebreak
The output I got when I ran this was:

\begin{center}
\line(1,0){400}
\end{center}
\footnotesize
\begin{Verbatim}[fontsize=\small,samepage=true,frame=none]
beard,height,lightsaber: 93.87
beard,demeanor,height,lightsaber: 93.40
beard,demeanor,lightsaber: 93.20
beard,lightsaber: 90.00
demeanor,height,lightsaber: 90.00
height,lightsaber: 90.00
demeanor,lightsaber: 90.00
lightsaber: 90.00
beard,demeanor,height: 85.80
beard,height: 85.60
\end{Verbatim}
\normalsize
\begin{center}
\line(1,0){400}
\end{center}

By a hair, the feature set that excludes \texttt{demeanor} but includes
everything else did the best, with 93.87\% accuracy. Right behind it was the
full feature set. (If we ran this multiple times, with different training set
choices, our results would vary somewhat.) This isn't too surprising, since in
this artificial data set, all of the features \textit{were} pretty indicative
of Jediness/Sithness, and they all \textit{were} fairly independent of each
other.

\subsubsection{Sadly...}

Now here's the deal. Although we might think we could declare victory with
this ``consider all the possible subsets and choose the best" technique, it
turns out that except in very small cases, \textit{it's absolutely
impossible.} Consider: what if instead of four features, we had forty?
(Certainly not a large number of features by \textit{any} means.) We would
have had:
\begin{align*}
2^{40} = 1,099,511,627,776
\end{align*}

or over a trillion different possible sets of features. Even if we could train
a \textit{thousand} different classifiers on our training data \textit{every
second}, here's how long it would take:

\footnotesize
\begin{align*}
1,099,511,627,776 \times 
    \frac{1\ \textrm{second}}{1000\ \textrm{classifiers}} \times
    \frac{1\ \textrm{minute}}{60\ \textrm{seconds}} \times
    \frac{1\ \textrm{hour}}{60\ \textrm{minutes}} \times
    \frac{1\ \textrm{day}}{24\ \textrm{hours}} \times
    \frac{1\ \textrm{year}}{365\ \textrm{days}} = 34.9\ \textrm{years}
\end{align*}

\normalsize
Ouch. You probably haven't even been alive for that long. And obviously the
situation gets worse the greater the number of features. If we had merely 100
features (also not an unreasonably large number at all), the universe will
have long since died of heat death by the time we're finished considering
subsets, even if we had every computer that had ever been made working
non-stop on the problem.

\subsection{``Greedy" approaches: forward \& backward selection}

So we have to find a way to choose a subset of features without considering
every different possibility. The most common way to do this is to use a
\textbf{greedy algorithm}, which means a procedure that finishes computing
quickly because it repeatedly takes (possibly non-optimal) shortcuts. More
precisely, a greedy algorithm always chooses \textit{the immediately
best-looking option} and goes with it, ignoring the fact that if it took more
time and considered further ahead, it would have made a different choice.

(To further illustrate a greedy algorithm: imagine a chess program that always
tried to take a piece if one was open to capture. It figures, ``hey, taking
the opponent's pieces is a good thing, and there's a piece sitting there
unprotected: why not take it?" Obviously such an algorithm might play well
against a newbie, but against even a modestly sophisticated player it would
suck, since it would be easily lured into traps. Without being able to
consider the longer-term consequences of what it's doing, it would make many
decisions that looked good in the near-term but which would soon be revealed
as dumb.)

\subsubsection{Forward selection}

The most common greedy algorithm used for feature selection is simply this:

\begin{itemize}
\itemsep.1em
\item Train $p$ different one-feature classifiers, where $p$ is the number of
features, and see how each one does. Choose as the winner the feature that
works the best in isolation.
\item Now, for each of the $p-1$ remaining features, train a classifier with
your original winner \textit{plus} the new one. Choose as the winner the
``best"\footnote{I put ``best" in quotes here because it may well not actually
be the best two-feature classifier out of all possible two-feature
classifiers, as we'll see.} \textit{two}-feature classifier.
\item Repeat this for each of the $p-2, p-3, p-4, \dots$ remaining features,
finding the best ``best" three-feature, four-feature, five-feature, $\dots$
classifier.
\item Your final answer is the best of the ``best." Out of the ``best"
one-feature, two-feature, \textit{etc.} classifiers, choose the one with the
best performance. This is the feature set you'll declare the winner.
\end{itemize}

Two things are important to see here. First, you're training and considering a
\textit{vastly} smaller number of classifiers than you are in the ``best
subset" approach we did previously. For the forty feature case, you'll be
considering:

\begin{align*}
\textrm{forward selection}&: 40 + 39 + 38 + \dots + 1 = 780\ \textrm{classifiers}\\
\textrm{best subset}&: 2^{40} = 1,099,511,627,776\ \textrm{classifiers}\\
\end{align*}

The time it will take to do this has gone from 35 years to .78
\textit{seconds}. Not bad.

The second thing to recognize is that we are \textit{not} by any means
guaranteed to get the best possible set of features this way. In fact, we
almost certainly won't. To see this, you have to realize that the feature that
looks the best of all in isolation might \textit{not} be one of the
\textit{pair} of features that works best as a pair. And this phenomenon
continues every iteration of our forward selection loop.

This can seem weird at first: if feature $F_{26}$ gives us the best
one-feature classifier of all forty, shouldn't feature $F_{26}$ plus another
one be the best \textit{two}-feature classifier? The answer is not
necessarily. Much of the reason why has to do with the overlap between
features -- recall that they will almost certainly contain redundant
information among them, and so if you knew from the outset that you wanted a
two-feature classifier, you might well choose $F_4$ and $F_{39}$ instead of
$F_{26}$, which was a great solo artist but didn't pair up as well with anyone
else.

\subsubsection{Backward selection}

As you might have guessed, you can do this entire process ``in reverse":

\begin{itemize}
\itemsep.1em
\item Start by including \textit{every} feature (instead of none of them).
\item Train $p$ different $p-1$-feature classifiers, each of which
\textit{leaves out} one feature. Choose as the winner the one with the
best performance.
\item Now, repeatedly leave out another feature, and choose as a winner the
best $p-2$-feature, $p-3$-feature, $p-4$-feature, $\dots$ classifier.
\item Your final answer is the best of the ``best." Out of the ``best"
$p-1$-feature, $p-2$-feature, \textit{etc.} classifiers, choose the one with
the best performance. This is the feature set you'll declare the winner.
\end{itemize}

Backward selection has the same advantages and disadvantages as forward
selection. Unsurprisingly, it will ordinarily \textit{not} produce the same
``best" feature set that forward selection does (nor that best-subset does).
Often data scientists perform both forward and backward selection and then
spend some time scrutinizing the feature sets each came up with to decide on
their final, best feature set.

\subsubsection{Evaluating classifier ``goodness"}

By the way, I repeatedly referred to ``best performance" in the description
above. What does this mean in practice? 

We know that more complex (``wiggly") classifiers will tend to overfit a
data set. Hence we can't blindly use performance on a training set to choose
between a four-feature classifier and a five-feature classifier, since the
five-feature one will \textit{almost certainly win for the wrong reasons.} If
it has more features to work with, it has more flexibility, and hence will
surely ``wiggle" closer to the points it trains on. This could be a blessing
but also a curse, depending on whether it's really picking up signal or
locking on to noise.

There are two different ways of dealing with this problem in order to
``fairly" compare classifiers of different complexities. I think of them as
the math-oriented and the computer-science-oriented ways:

\begin{itemize}
\itemsep.1em

\item \textbf{The ``math-oriented" way.} If you use the entire data set to
train on, you can still adjust for the complexity factor by
\textit{penalizing} more complex models in a principled way. The details are
outside the scope of this class, but I'll throw the following buzzterms at you
so that if you see them you'll know what category to file them under:
w that more complex (``wiggly")
classifiers will tend to overfit a training set. Hence we can't blindly use
the training-set 

\textbf{``Mallows' $C_p$"}, \textbf{AIC}, \textbf{BIC}, and \textbf{adjusted
$R^2$}. All of these are ways of taking a model's error rate and adjusting it
in light of how complex that model is, and therefore how likely it is to
overfit.

\item \textbf{The ``computer-science-oriented" way.} Just use a separate
training set and test set. Drr.

\end{itemize}

\small
\subsection{Shrinkage}

For completeness, I briefly mention two other common methods of feature
reduction. The first, most commonly used in linear regression/classification,
goes by the general term ``\textbf{shrinkage}," and its two most common
variants are called \textbf{ridge regression} and \textbf{lasso regression}.
These can both be thought of as ways of ``partially including" features. Each
of them has a parameter you can dial to control how aggressively they will
``shrink" coefficients to zero. This helps when some of your features are
correlated with others, since dialing down their influence will prevent them
from jointly dominating the prediction. The mathematics are subtle, although
one point worth mentioning is that with lasso, feature coefficients can
actually be reduced all the way to \textit{zero}, which effectively eliminates
them and performs a kind of feature selection as in the previous section.
(Ridge regression does not have this property.)

\subsection{Transforming the feature space}

Finally, one of the coolest approaches to reducing the feature space is to
\textit{transform it into a different set of dimensions}. This may sound
heady, and it is: it involves some linear algebra which you may know or learn
someday. But the basic idea is that the raw features themselves, as given in
the data set, may not be the most useful measurements of the data: perhaps
some \textit{combinations} of the features can give us more mileage. Instead
of height, age, salary, and work experience, perhaps we can take some
mathematical function of height-and-age, and another mathematical function of
salary-and-work-experience, and work with those two combined features instead.
(Often, the combined features will involve components of \textit{all} the
original features, not just in pairs as in this example.) The general term
\textbf{dimensionality reduction} is what to look for in this area, together
with specific techniques like \textbf{principal components analysis} (PCA) and
\textbf{singular value decomposition} (SVD).


